\subsection{The actual root-point Fourier projection}
\label{subsec:ftb-root-fourier-projection}

We retain one rational Hermitian space \(X\), its symplectic
coordinate map \(\mathrm{bas}\), and its genuine Heisenberg action
\(U_y=\mathrm{heisG}(X,\mathrm{bas},y)\). All lattices in this
subsection are subgroups of this one space, and all test functions
belong to its one Schwartz--Bruhat representation.
Put
\[
 \Omega(y,z)=\sum_{b<n}EF_X(y_b,z_b),\qquad
 \epsilon(y)=\mathrm{epsW}(X,w,y),\qquad
 \rho(y)=\epsilon(y)U_y .
\]
The actual action has
\[
 U_yU_z=e^{-\pi i\Omega(y,z)}U_{y+z},\qquad
 U_yU_z=e^{-2\pi i\Omega(y,z)}U_zU_y.             \tag{RP1}
\]
The values of \(EF_X\) and hence \(\Omega\) are rational.
Take \(n\) even and \(w\in\mathcal O_E\) with \(w+\bar w=1\).

Here are the exact lattice and marking hypotheses used below.
Let \(l\) be the source level, write
\[
 P=\operatorname{Mfr}(l),\quad Q=l.\operatorname{Mb}=\bar P,
 \quad M=l.M,\quad P+Q=\mathcal O_E,\quad PQ=(M).
\]
In particular \(M\geq3\) is odd. Let \(L\subset E^r\) be finitely
generated and selfdual for \(EF_X\), in the literal sense
\[
 x\in L\quad\Longleftrightarrow\quad
                         EF_X(x,L)\subset\mathbb Z.       \tag{RP2}
\]
Let \(\Lambda\in\operatorname{Vset}(E,n,M)\): it is a finitely
generated selfdual lattice for the standard Hermitian form, and
there is \(s>0\), prime to \(M\), with
\[
 s\Lambda\subset\mathcal O_E^n,\qquad
                         s\mathcal O_E^n\subset\Lambda.   \tag{RP3}
\]
The number \(s\) is an auxiliary denominator, unrelated to a
choice of roots. Write \(T(\Lambda,B)\) for the additive group
generated by the matrices \(v_bu_j\), with \(v\in\Lambda,u\in B\).
A usable level \(N\) means precisely
\[
 N\text{ odd},\quad M\mid N,\quad
                N\mathcal O_E^n\subset\Lambda\cap\mathcal O_E^n.
                                                               \tag{RP4}
\]

The abstract level lattice, a finite \(\mathcal O_E\)-module,
is denoted by \(\mathcal L\).
The marking \(m:\mathcal L\to E^r\) is \(\mathcal O_E\)-linear and
has exactly the following three properties:
\[
 \begin{split}
 &m(\mathcal L)\subset L;\\
 &a\in P\mathcal L,\ m(a)\in ML
                              \ \Longrightarrow\ a\in M\mathcal L;\\
 &v\in PL\ \Longrightarrow\
           v=m(a)+Mu\quad(a\in P\mathcal L,\ u\in L).
 \end{split}                                                   \tag{RP5}
\]
These are \(\mathrm{IsMark}\), rather than an unrecorded assertion
that \(m\) is an isomorphism. Put
\(\mathcal Q=(\mathcal L/Q\mathcal L)^n\).
The supplied root representatives \(Y_q\in\operatorname{Mat}_{n,r}(E)\)
satisfy, for every \(q\) and row \(b\),
\[
 \begin{split}
 &a\in P\mathcal L,\quad [a]=q_b,\quad
                         Y_{q,b}-M^{-1}m(a)\in L
                         \quad\hbox{for some }a;\\
 &Q^jY_{q,b}\subset 2NL\quad\hbox{for some integer }j\geq0.
 \end{split}                                                   \tag{RP6}
\]
The exponent \(j\) may depend on \(q,b\). This is the original
\(\mathrm{IsRoots}\); it neither replaces roots by arbitrary lifts
nor asserts compatibility of different choices of \(N\).

\begin{srlemma}[Tensor contraction and the ideal shift]
\label{lem:ftb-root-tensor-shift}
Define
\[
 L^\sharp=(PL:Q)=\{x:Qx\subset PL\},\qquad
 A=T(\Lambda,L),\quad A^\sharp=T(\Lambda,L^\sharp).
\]
Then \(L^\sharp\) is finitely generated and selfdual for \(EF_X\).
For either \(B=L\) or \(L^\sharp\),
\[
 y\in T(\Lambda,B)
 \quad\Longleftrightarrow\quad
 C_y(v):=\sum_b\overline{v_b}y_b\in B
                                      \quad(v\in\Lambda). \tag{RP7}
\]
Both \(A,A^\sharp\) are the source's good signed selfdual lattices.
In particular \(\rho\) is additive on each, their signed vacuum
line is one dimensional, the normalized vacuum has norm one,
and its translate by a vector outside that lattice is orthogonal
to it. Also \(\epsilon(2a)=1\) for \(a\) in either lattice.
\end{srlemma}

\begin{proof}
For (RP7), a finite integral dual family \(v_i,w_i\in\Lambda\)
satisfies
\[
 z=\sum_i\langle z,w_i\rangle v_i,\qquad
                      y=\sum_i v_i\otimes C_y(w_i).
\]
Such a family is obtained by splitting the finite projective
\(\mathcal O_E\)-module \(\Lambda\) off a finite free module and
using its Hermitian selfduality to represent the dual coordinate
maps. Contraction of a pure tensor is
\(C_{v\otimes u}(w)=\langle v,w\rangle u\).
This proves both directions of (RP7), with no freeness assumption
on either \(\mathcal O_E\)-lattice.

We spell out the source's ideal-shift argument. Choose
\(q_0\in Q,p_0\in P\) with \(q_0+p_0=1\). Multiplication by \(M\)
embeds \(L^\sharp\) into \(L\), since \(M\in Q\) and
\(ML^\sharp\subset PL\subset L\). Thus finite generation follows
from the Noetherian property of \(\mathcal O_E\).
The adjoint identity is
\[
                 EF_X(au,v)=EF_X(u,\bar a v).              \tag{RP8}
\]
It implies
\[
                        EF_X(PL,L^\sharp)\subset\mathbb Z:
\]
on a generator \(a\ell\) of \(PL\), move \(a\) to the second
slot; \(\bar a\in Q\), so \(\bar a y\in PL\subset L\), and
apply (RP2). If \(x,y\in L^\sharp\), split \(x=q_0x+p_0x\).
The first term pairs integrally with \(y\). Move \(p_0\) in the
second term to \(\bar p_0y\in PL\), then use alternation.
This proves integrality on \(L^\sharp\).

Conversely suppose \(EF_X(x,L^\sharp)\subset\mathbb Z\).
For \(\nu\in Q\) and \(\ell\in L\), the vector
\(\bar\nu\ell\) belongs to \(L^\sharp\), since
\(Q\bar\nu\ell\subset PL\). Equation (RP8) and (RP2) give
\(\nu x\in L\). For \(\nu_1,\nu_2\in Q\), the vector
\[
                M^{-1}\bar\nu_1\bar\nu_2\ell
\]
also belongs to \(L^\sharp\): multiplying it by \(b\in Q\),
the factor \(b\bar\nu_1\) is in \(QP=(M)\), leaving a factor
\(\bar\nu_2\in P\). Hence
\(M^{-1}\nu_1\nu_2x\in L\), and
\(\nu_1\nu_2x\in ML\subset PL\). Finally, for \(b\in Q\),
\[
                   bx=(q_0b)x+p_0(bx)\in PL.
\]
Thus \(x\in L^\sharp\), proving selfduality.

For clarity, the last signed-vacuum assertions use the same
tensor-lattice vacuum construction as the surrounding theta
system, with these literal \(X,\mathrm{bas},w,n\).
Tensor selfduality follows from (RP7) and (RP2), since pairing
against pure tensors is \(EF_X(C_y(v),u)\). The semicharacter
identity makes \(\rho\) additive on an integral tensor lattice;
its square identity gives \(\epsilon(2a)=1\).
In the rational symplectic coordinates \(\mathrm{bas}\), the
finite Heisenberg projection constructs the unit vacuum and
its unique fixed line. Orthogonality has the following direct
verification used in that construction. If \(y\notin A\),
selfduality gives \(a\in A\) with \(\Omega(y,a)\notin\mathbb Z\).
The commutator (RP1) and unitarity make
\[
 \langle\rho(y)v,v\rangle
          =e^{2\pi i\Omega(y,a)}\langle\rho(y)v,v\rangle=0.
\]
The same statement holds with \(A^\sharp\).
Thus these conclusions use the actual signed selfdual-vacuum
producer; a bare claim about a classical connected moduli space
is not a replacement for it.
\end{proof}

\begin{srlemma}[The actual halves and quotient representatives]
\label{lem:ftb-root-halves}
For all the data above,
\[
 Y_q/2\in A^\sharp,\quad
       \epsilon(Y_q)=1,\quad
 u\in T(\Lambda,2QL)\Longrightarrow u/2\in A,\
                                      \epsilon(u)=1.       \tag{RP9}
\]
Moreover
\[
 \begin{split}
 &Y_q-Y_{q'}\in A\quad\Longrightarrow\quad q=q',\\
 &y\in A^\sharp\quad\Longrightarrow\
                      \exists q,\ y-Y_q\in A .
 \end{split}                                              \tag{RP10}
\]
\end{srlemma}

\begin{proof}
First (RP5)--(RP6) imply \(QY_{q,b}\subset L\):
if \(a\in P\mathcal L\) is the representative in (RP6) and
\(t\in Q\), then \(ta\in QP\mathcal L=M\mathcal L\), so
\(tM^{-1}m(a)\in L\).

We give the denominator calculation needed to prove the first
assertion of (RP9). Fix \(q\), choose \(s\) as in (RP3), and take
\(J\) at least the sum of the row exponents in (RP6).
For \(v\in\Lambda\) put \(z=C_{Y_q/2}(v)\).
For \(t\in Q^J\), write \(tY_{q,b}=2Nu_b\), \(u_b\in L\).
Then
\[
                         tz=\sum_b\bar v_bNu_b\in PL.     \tag{RP11}
\]
Indeed choose integers \(a,b'\) with \(as+b'M=1\).
Selfduality of \(\Lambda\) and \(N\mathcal O_E^n\subset\Lambda\)
imply \(N\Lambda\subset\mathcal O_E^n\).
Thus both \(sv_b,Nv_b\) are integral, and each summand decomposes as
\[
 \bar v_bNu_b
    =aN\,\overline{sv_b}u_b
                   +b'M\,\overline{Nv_b}u_b\in PL,
\]
because \(M,N\in P\).

For \(c\in Q\), we also have
\[
                          2s\,cz\in L:                  \tag{RP12}
\]
it is the contraction of the integral rows \(cY_{q,b}\in L\)
against \(sv\in\mathcal O_E^n\).
Since \(Q^J+P=\mathcal O_E\), combine (RP11) for \(cz\)
with (RP12) to obtain \(2s\,cz\in PL\).
Since \(M\) is odd and \(s\) is prime to \(M\),
\(\gcd(2s,M^J)=1\); another integer Bezout combination, using
\(M^J\in Q^J\), gives \(cz\in PL\).
Therefore \(z\in L^\sharp\), and (RP7) proves \(Y_q/2\in A^\sharp\).
The assertion \(u/2\in A\) follows by contracting
\(u\in T(\Lambda,2QL)\) and dividing its factor \(2\).
The signs now follow from \(\epsilon(2a)=1\).

We record two row calculations, since they prevent a replacement
of the original quotient by an assumed coordinatewise quotient.
If \(z\in T(\Lambda,L)\) and \(Mz_b\in L\), then \(sz_b\in L\)
by (RP7) applied to \(se_b\in\Lambda\); (RP3) and
\(\gcd(s,M)=1\) give \(z_b\in L\).
Next if \(y\in A^\sharp\), set \(y'=as\,y\), with
\(as+b'M=1\). Then \(My'_b\in PL\).
For any \(Y\in A^\sharp\) with \(Y_b-y'_b\in L\) for every row,
\[
                              y-Y\in A.                 \tag{RP13}
\]
To verify this, contract \(z=y-Y\) against \(v\in\Lambda\).
This contraction lies in \(L^\sharp\), so its \(M\)-multiple
lies in \(L\). Its \(s\)-multiple also lies in \(L\), by
\[
 y-Y=b'My-(Y-as\,y)
\]
and \(sv\in\mathcal O_E^n\). Bezout in \(s,M\) makes the
contraction itself integral, proving (RP13).

If \(Y_q-Y_{q'}\in A\), (RP6) shows its \(M\)-multiple has
every row in \(L\); the preceding row calculation gives
\(Y_{q,b}-Y_{q',b}\in L\).
For representatives \(a,a'\in P\mathcal L\), this implies
\(m(a-a')\in ML\). The injectivity clause of (RP5) gives
\(a-a'\in M\mathcal L\subset Q\mathcal L\), hence \(q=q'\).

Finally take \(y\in A^\sharp\) and its \(y'\) from (RP13).
The surjectivity clause of (RP5) writes
\(My'_b=m(a_b)+Mv_b\), with \(a_b\in P\mathcal L,v_b\in L\).
Choose \(q_b=[a_b]\). The representative \(a'_b\) used for this
root lies in the same \(Q\)-class, and
\[
 a'_b-a_b\in P\mathcal L\cap Q\mathcal L
                            =PQ\mathcal L=M\mathcal L.
\]
This equality follows by a Bezout combination of \(P,Q\);
it needs no freeness of \(\mathcal L\).
Thus \(Y_{q,b}-y'_b\in L\). Since \(Y_q\in A^\sharp\),
(RP13) proves \(y-Y_q\in A\).
\end{proof}

For a good signed selfdual lattice \(B\), the source projection
\(\operatorname{proj}_B p\) is the average of the \emph{distinct}
vectors in the finite signed orbit
\(\{\rho(y)p:y\in B\}\). Its value at \(p=0\) is zero.
In particular one must identify these distinct vectors before
replacing this average by a sum over \(\mathcal Q\).

\begin{srproposition}[The original root Fourier identity]
\label{prop:ftb-actual-root-fourier}
Let \(p\) be Schwartz--Bruhat with \(\rho(a)p=p\) for \(a\in A\).
For every \(u\in T(\Lambda,2QL)\),
\[
 \sum_{q\in\mathcal Q}e^{-2\pi i\Omega(u,Y_q)}U_{Y_q}p
       =|\mathcal Q|\,U_u\bigl(\operatorname{proj}_{A^\sharp}p\bigr).
                                                               \tag{RP14}
\]
The action, signs, root representatives, shifted lattice and
projection are all the literal ones fixed above.
\end{srproposition}

\begin{proof}
If \(p=0\), both sides vanish. Assume \(p\ne0\).
The unique vacuum line and the preceding orthogonality show that
the full signed stabilizer of \(p\) is exactly \(A\).
Indeed \(p=cv_A\), \(c\ne0\); if \(\rho(y)p=p\) with \(y\notin A\),
its inner product with \(p\) would be both zero and
\(\|p\|^2>0\).

Every \(y\in A^\sharp\) differs from some \(Y_q\) by a vector
of \(A\), by (RP10), and that difference is also in \(A^\sharp\).
Additivity of \(\rho\) on \(A^\sharp\) therefore gives
\(\rho(y)p=\rho(Y_q)p=U_{Y_q}p\), using (RP9).
If \(U_{Y_q}p=U_{Y_{q'}}p\), multiply by
\(\rho(-Y_{q'})\) inside this same representation.
The result is \(\rho(Y_q-Y_{q'})p=p\), so the exact stabilizer
and (RP10) give \(q=q'\).
Consequently the signed orbit consists of the distinct vectors
\(U_{Y_q}p\), and
\[
                 \operatorname{proj}_{A^\sharp}p
                         ={1\over|\mathcal Q|}
                                     \sum_q U_{Y_q}p.     \tag{RP15}
\]
The quotient is finite: \(\mathcal L\) is a finite
\(\mathcal O_E\)-module and \(Q\ne0\); it is also nonempty.
For the actual \(u\), (RP9) and the fixed-vector condition give
\(U_up=p\). Now (RP1) gives, term by term,
\[
 U_uU_{Y_q}p=e^{-2\pi i\Omega(u,Y_q)}U_{Y_q}U_up
                     =e^{-2\pi i\Omega(u,Y_q)}U_{Y_q}p .
\]
Apply \(U_u\) to (RP15) and multiply by \(|\mathcal Q|\).
\end{proof}

\paragraph{Exact frozen-source correspondence and remaining inputs.}
All paths in this paragraph are at
\srcid{alpoge-source-7cbf7a0c}.
The marking and root conditions (RP5)--(RP6), including the
\(Q\)-power condition with \(2N\), are
\srcid{src/Development/AbelianVarietyTheory/Definitions/SpanningFamilies/IsStabForm.lean},
(\srcid{ca69affd0453770815104e7bfc7b4a93c7ef988e}), declarations \srcid{IsMark}, \srcid{IsRoots}.
The level \(N\) is \srcid{GenLat.LevOK} in
\srcid{src/Formalism/LocalToGlobal/ProjectiveLattices/ThShort.lean}
(\srcid{ba33f03086ab6eb7e86ad7e31fe5d9484953af5e}).
The tensor criterion is \srcid{TensLat.mem\_tens\_iff} in
\srcid{src/Development/AbelianVarietyTheory/Proofs/SiegelThetaTransformation/TensorLatticeSelfDualSemicharacterIndex.lean}
(\srcid{4894f884d42f8c17e88ea2604bcf2fc40817b20d});
the signed vacuum, orthogonality and actual finite-orbit projection are
\srcid{TensVac.Good.vac}, \srcid{Good.ip\_eq\_zero},
\srcid{Good.finite} in
\srcid{TensorLatticeVacuaExistUnique.lean}
(\srcid{226ea09aeee65b72e52c8447abc4d4f010b3db15}).
The original endpoint is \srcid{GlobalTheory.Dict.poisson\_roots} in
\srcid{FourierTranslatesVacuumRootsLevelPoints.lean}
(\srcid{28fc115f8c46ce96e6f2bbf704fe6c4b1b43aced}).
Its \srcid{selfD\_shift}, \srcid{mem\_of\_two},
\srcid{half\_away}, \srcid{half\_mem}, \srcid{rows\_of\_mem},
\srcid{exists\_rows}, \srcid{mem\_of\_fix} and
\srcid{fourier\_of\_reps} are precisely the intermediate arguments
above. No large-rank inequality is used here; the statement keeps
all even \(n\), including \(n=0\).
This closes the root-point projection identity retained as (RS9)
in the short-reading comparison, from the displayed actual
marking, roots and signed-vacuum inputs. It does not construct
the level cover, its roots, or the common algebraic/analytic
marking; those are separate suppliers.
